\subsection{Necesidad de seleccionar un simulador base}

La propuesta de este trabajo no se limita a ejecutar circuitos cuánticos sobre una
herramienta existente, sino que exige intervenir la representación del vector de estado
para integrar una estrategia de compresión sin pérdida de precisión y evaluar su efecto
sobre memoria, tiempo y exactitud numérica. Por esta razón, la selección del simulador
base constituye una decisión metodológica central, ya que condiciona la viabilidad de
la integración, el alcance de la instrumentación experimental y el esfuerzo técnico
requerido para desarrollar la solución.

En los simuladores cuánticos tipo Schrödinger, el vector de estado es la estructura
principal del cálculo y, a la vez, la fuente dominante del consumo de memoria. En
consecuencia, el simulador seleccionado debía permitir una intervención suficientemente
directa sobre esa representación, sin desplazar el problema hacia restricciones
derivadas de abstracciones internas difíciles de modificar o de mecanismos poco
accesibles para observación experimental.

Con este propósito se consideraron cinco alternativas representativas dentro del
ecosistema de simulación cuántica de estado completo: QuEST, Intel-QS, qsim,
Quantum++ y TMFQSfullstate. La elección se formalizó mediante criterios explícitos y
una matriz de evaluación ponderada, de modo que la decisión quedara vinculada a los
objetivos técnicos y experimentales del trabajo.
